\section{Related Work}

\subsection{Software Engineering Benchmarks}
Code generation benchmarks, which evaluate models on the task of synthesizing code from natural language descriptions, have served as a long-standing bellwether for measuring LM performance~\citep{chen2021evaluating,austin2021program,hendrycks2021measuring,lu2021codexglue}. 
Subsequent works have built upon the code generation task formulation to contribute new benchmarks that translate problems to different (programming) languages~\citep{cassano2022multiple,wang2023mconala}, incorporate third-party libraries~\citep{lai2022ds1000,liu2024mlbench}, introduce derivative code completion tasks~\citep{huang2024mlagentbench,muennighoff2023octopack}, increase test coverage~\citep{evalplus}, change the edit scope~\citep{ding2023crosscodeeval,du2023classeval,Yu2023CoderEvalAB}, and add robustness to dataset contamination~\citep{jain2024livecodebench}.
Code generation problems are largely self-contained, with short problem descriptions ($\sim$100 lines) and corresponding solutions that are similarly brief, requiring nothing more complex than basic language primitives.
Tests are either handwritten or generated synthetically via fuzz testing.
In recent months, the rapid development of LMs has begun to saturate many of these benchmarks.
For instance, the top method solves \num{94.4}\% of HumanEval~\citep{zhou2023language}.

Gauging future trends with the code generation task paradigm can be limited by the simplicity of this setting and cost of human-in-the-loop problem creation.
In response, recent efforts have demonstrated that software engineering (SE) can serve as a diverse, challenging testbed for LM evaluation~\citep{zhang2023repocoder,jimenez2024swebench,Liu2023RepoBenchBR}.
Repository-level code editing introduces many reasoning challenges grounded in real SE subtasks, such as spotting errant code and identifying cross-file relationships and understanding codebase-specific symbols and conventions.
As a field, SE has generally studied tasks in a more isolated manner; prior benchmarks tended to frame problems in isolation from the rest of a codebase~\citep{justJE2014defects,Karampatsis2019HowOD}.

We use SWE-bench because it unites many separate SE tasks, such as automated program repair~\citep{fan2023automated,sobania2023analysis,xia2022less}, bug localization~\citep{chakraborty2018entropy,yang2024localization}, and testing~\citep{kang2023large,wang2023software,xia2023universal} under a single task formulation that faithfully mirrors practical SE.
Furthermore, SWE-bench task instances are diverse, having been automatically collected from real GitHub issues across \num{12} different repositories.
In addition, SWE-bench performance is based on rigorous, execution-based evaluation with human-written unit tests.

\subsection{Language Models as Agents}
The co-emergence of stronger LMs, increasingly challenging benchmarks, and practical use cases have together motivated a paradigm shift in LMs' inference setting.
Instead of traditional zero/few-shot generation, LM agents~\citep{hong2023metagpt,sumers2023cognitive,wang2024survey,xi2023rise} that interact with a real/virtual world have proliferated as the default setting for web navigation~\citep{koh2024visualwebarena,nakano2022webgpt,press-etal-2023-measuring,sridhar2023hierarchical,thoppilan2022lamda,yao2023webshop,yao2023react,zhou2023webarena}, computer control~\citep{packer2024memgpt,wu2024oscopilot,xie2024osworld}, and code generation tasks~\citep{holt2024lmac,wang2023executionbased,yin2022natural}.


Interaction and code generation are increasingly used together, with code as the modality of choice for actions~\citep{wang2024executable,yang2023intercode}, tool construction~\citep{gu2024middleware,wang2024trove,zhang2024training}, and reasoning~\citep{shinn2023reflexion,zelikman2022parsel,zelikman2024selftaught}.
Coding agents have also been applied to offensive security~\citep{fang2024llm,shao2024empirical,yang2023language}, theorem proving~\citep{thakur2024incontext}, and clinical tasks~\citep{shi2024ehragent,tang2024medagents,wornow2024automating}.
To the best of our knowledge, SWE-agent is the first work to explore language agents for end-to-end software engineering (SE).
% SWE-agent's results reliably indicate the robustness and scalability of agents' ability to reason with code.
